\section{Discussion}

The pursuit of robust multi-agent systems (MAS) in the face of non-stationarity has traditionally been hampered by the ``moving target'' problem, where the co-evolution of agent policies invalidates the Markov assumption \cite{Foerster2017}. Our work addresses this by reframing communication not as a secondary signal to be optimized via reward, but as a primary mechanism for Bayesian state-matching within an active inference framework.

\subsection{Synthesis of Results and Theoretical Implications}
Our experiments demonstrate that by minimizing the Expected Free Energy (EFE), agents naturally adopt communication strategies that prioritize resolving uncertainty about their neighbors' intentions and the environmental dynamics. Unlike standard MARL approaches that often converge to brittle, over-fitted communication protocols \cite{Sukhbaatar2016}, our active inference formulation maintains a generative model that explicitly accounts for the epistemic value of communication. This aligns with recent theoretical insights into ``Free-Energy Equilibria'' \cite{ref_new_FreeEnergyEquil}, suggesting that boundedly-rational agents can achieve stable coordination in non-stationary regimes by minimizing the divergence between their internal models and the empirical distribution of messages.

Furthermore, our results indicate that agents utilizing active inference are more resilient to the ``missing reward'' problem \cite{ref_new_TheMissingRewar}. While reinforcement learning relies on dense, stationary reward signals to guide communication, our agents rely on the intrinsic drive to minimize prediction error. This intrinsic motivation allows for rapid adaptation when environmental transitions shift unexpectedly, as the communication policy is coupled to the agent's updated belief about the world's hidden states rather than a historical trajectory of rewards.

\subsection{Non-Stationarity as a Latent Challenge}
A critical observation in our study is that non-stationarity acts as a source of ``dynamic correlation'' between agents. Conventional methods often struggle to adjust to these shifts in a timely manner. Our approach, mirroring findings in active inference-driven multi-armed bandits \cite{ref_new_ActiveInference}, allows for dynamic correlation adjustments. By treating the policies of other agents as hidden variables to be inferred through communication, our framework effectively transforms a non-stationary problem into a stationary one within a higher-order belief space. This shift from ``policy learning'' to ``belief tracking'' is what enables the superior performance observed in our non-stationary predator-prey and coordination benchmarks.

\subsection{Limitations and Scalability}
Despite the principled nature of our framework, several limitations persist. First, the computational complexity of calculating the EFE, particularly the information gain term, scales poorly with the number of agents and the complexity of the message space. While we employed variational approximations to maintain tractability, true scalability to massive MAS remains an open challenge. Second, our current model assumes a degree of shared prior knowledge regarding the communication protocol's structure. In fully heterogeneous settings where agents possess entirely different generative models, the ``Bayesian blurring'' effect may impede the alignment of internal states.

\subsection{Future Horizons and Emerging Paradigms}
The intersection of active inference and large-scale computing offers fertile ground for future research. One promising direction is the integration of our framework with Large Language Models (LLMs) for high-level reasoning and resource allocation. As suggested by recent work on active inference for cloud-edge computing \cite{ref_new_Largelanguagemo}, LLMs could serve as the ``generative core'' for MAS, providing the semantic richness necessary for more complex communication acts in decentralized environments.

Additionally, exploring incentive design within the active inference framework \cite{ref_new_ActiveInference} could provide a path toward aligning competitive agents in non-stationary games. By engineering the ``free energy landscape'' of a system, a central designer could encourage emergent communication that leads to socially optimal outcomes without requiring explicit control over individual agent policies.

\section{Conclusion}
In this paper, we presented a principled framework for multi-agent communication grounded in active inference. By treating communication as an active process of minimizing variational free energy, we have shown that agents can effectively navigate the challenges of non-stationarity, achieving superior coordination and adaptation compared to state-of-the-art MARL baselines. Our work bridges the gap between cognitive science-inspired inference and multi-agent engineering, providing a robust foundation for future autonomous systems that must communicate to survive and thrive in an ever-changing world.